% ------------------------------------------------------------------
% ch-results-synth.tex --- body of the chapter
%   "Results on Synthetic Data"  (\chapter in main.tex)
% Do not add a \chapter command here.
% ------------------------------------------------------------------

This chapter is intentionally empty of numbers. It is the destination for the protocols of
\cref{ch:codeplan}: every synthetic-data run writes its verdict here, append-only, under the artifact
rules of \cref{chcp:sec:ci}. Each section below mirrors one protocol of \cref{chcp:sec:forward}. A
section is filled only after its protocol has run under its frozen pass criterion; until then its
table shell stays empty. Real-data results never appear here; they belong to \cref{ch:results-spain}.

\section{Generator sanity and stability checks}\label{chrs:sec:generator}

Before any fitter is judged, the rehearsal worlds themselves must pass: stable event counts across
seeds, no supercritical excursions, no runaway per-firm deal concentration, and a deduplication rule
that keeps legitimate repeat fundings. These checks apply the harness rules of \cref{chcp:sec:harness}
to the generator itself.

\begin{table}[htbp]
\centering\small
\begin{tabular}{lccc}
\toprule
Check & Regime A & Regime B & Verdict \\
\midrule
\multicolumn{4}{l}{\todo{fill after running the generator checks of \cref{chcp:sec:harness}}} \\
\bottomrule
\end{tabular}
\caption{Generator sanity and stability checks, per rehearsal regime.}
\label{chrs:tab:generator}
\end{table}

\section{Poisson baselines and the covariates-vs-excitation decomposition}\label{chrs:sec:b2}

Results of protocol B2 (\cref{chcp:prot:b2}): the nested NLL ladder from seasonal-naive through
covariates-only Poisson to covariates-plus-excitation, decomposed rung by rung with error bars.

\begin{table}[htbp]
\centering\small
\begin{tabular}{lccc}
\toprule
Ladder rung & Held-out NLL (A) & Held-out NLL (B) & $\Delta$ vs.\ previous rung \\
\midrule
\multicolumn{4}{l}{\todo{fill after running protocol B2 (\cref{chcp:prot:b2})}} \\
\bottomrule
\end{tabular}
\caption{Improvement decomposition across the model ladder.}
\label{chrs:tab:b2}
\end{table}

\section{Parametric Hawkes recovery and bucket-width identifiability}\label{chrs:sec:e1}

Results of protocol E1 (\cref{chcp:prot:e1}): recovery of $\branch$ and $\decay$ at 1d, 7d, and 28d
bucket widths, and the fitted-radius curve against latent-factor strength.

\begin{table}[htbp]
\centering\small
\begin{tabular}{lcccc}
\toprule
Bucket width & $\branch$ bias (SE) & $\decay$ bias (SE) & Ridge visible & Verdict \\
\midrule
\multicolumn{5}{l}{\todo{fill after running protocol E1 (\cref{chcp:prot:e1})}} \\
\bottomrule
\end{tabular}
\caption{Bucket-width identifiability and latent-confounding results.}
\label{chrs:tab:e1}
\end{table}

\section{Nonparametric challenger comparison}\label{chrs:sec:nonparam}

Results of protocol B3 (\cref{chcp:prot:b3}): the intensity models of \cref{ch:nonparam} against the
log-linear GLM and the parametric Hawkes, on held-out likelihood and calibration, in the same harness.

\begin{table}[htbp]
\centering\small
\begin{tabular}{lccc}
\toprule
Model & Held-out NLL & Calibration (PIT / dispersion) & Verdict \\
\midrule
\multicolumn{4}{l}{\todo{fill after running protocol B3 (\cref{chcp:prot:b3})}} \\
\bottomrule
\end{tabular}
\caption{Nonparametric challenger versus parametric incumbents.}
\label{chrs:tab:nonparam}
\end{table}

\section{Choice-model and ranking metrics}\label{chrs:sec:ranking}

Results of protocol E2 (\cref{chcp:prot:e2}) and benchmark B1 (\cref{chcp:prot:b1}): choice-model
design choices, weight recovery, and ranking quality under the recall-at-$K$ protocol of
\cref{ch10:sec:protocol}.

\begin{table}[htbp]
\centering\small
\begin{tabular}{lcccc}
\toprule
Model & Held-out NLL & Recall@5 & Recall@10 & MRR \\
\midrule
\multicolumn{5}{l}{\todo{fill after running protocols E2 (\cref{chcp:prot:e2}) and B1 (\cref{chcp:prot:b1})}} \\
\bottomrule
\end{tabular}
\caption{Choice-model and ranking results under the protocol of \cref{ch10:sec:protocol}.}
\label{chrs:tab:ranking}
\end{table}

\section{Open-universe and newcomer calibration}\label{chrs:sec:newcomer}

Results of the newcomer-design sweep (E2.2 of \cref{chcp:prot:e2}) and the outside-option calibration
targets of \cref{ch:universe}: quarterly newcomer-share calibration and coverage diagnostics.

\begin{table}[htbp]
\centering\small
\begin{tabular}{lccc}
\toprule
Design & Newcomer-share error (pp) & Incumbent weights stable & Verdict \\
\midrule
\multicolumn{4}{l}{\todo{fill after running E2.2 of \cref{chcp:prot:e2}}} \\
\bottomrule
\end{tabular}
\caption{Newcomer designs and open-universe calibration.}
\label{chrs:tab:newcomer}
\end{table}

\section{Causal-pipeline validation}\label{chrs:sec:causal}

Results of the causal validation protocol (\cref{chcp:prot:causal}): naive, backdoor, and AIPW
estimates against the known estimand, with the four coded assertions.

\begin{table}[htbp]
\centering\small
\begin{tabular}{lccc}
\toprule
Estimator & Estimate (CI) & Gap to truth & Assertion \\
\midrule
\multicolumn{4}{l}{\todo{fill after running the causal validation protocol (\cref{chcp:prot:causal})}} \\
\bottomrule
\end{tabular}
\caption{Causal-pipeline validation on known ground truth.}
\label{chrs:tab:causal}
\end{table}

\section{Walk-forward backtest}\label{chrs:sec:e5}

Results of protocol E5 (\cref{chcp:prot:e5}): the end-to-end weekly walk-forward, scored against
oracle and random, with the covariate ablations.

\begin{table}[htbp]
\centering\small
\begin{tabular}{lcccc}
\toprule
Variant & Top-5 & MRR & PIT calibration & Lead time (weeks) \\
\midrule
\multicolumn{5}{l}{\todo{fill after running protocol E5 (\cref{chcp:prot:e5})}} \\
\bottomrule
\end{tabular}
\caption{Walk-forward backtest with covariate ablation.}
\label{chrs:tab:e5}
\end{table}
